\documentclass[11pt]{article}
\usepackage[margin=.75in,top=.6in,bottom=.6in]{geometry}
\usepackage{amsmath,amssymb,amsthm,hyperref}
\renewcommand{\normalsize}{\fontsize{10}{12}\selectfont}
\DeclareMathSizes{10}{10}{8}{6}
\hypersetup{colorlinks=true,linkcolor=blue,urlcolor=blue,citecolor=blue}
\setlength{\emergencystretch}{2em}
\newtheorem{theorem}{Theorem}
\newtheorem{lemma}{Lemma}
\newcommand{\per}{\operatorname{per}}
\begin{document}
\begin{center}
{\Large\bfseries Sharp constants for a permanental\\Koteljanskii ratio in dimension three}\\[.5em]
Research proof and Lean formalization\\
October 26, 2026
\end{center}
\begin{abstract}
For every real totally positive $3\times3$ matrix $A$, we prove
\[
 \frac23<\frac{\per(A_{12,12})\per(A_{23,23})}{\per(A)a_{22}}<\frac98.
\]
Both constants are optimal within the strictly totally positive class and
are not attained there. Positive row and column scaling reduces the proof
to two explicit polynomial positivity certificates in four variables.
Strictly totally positive one-parameter families approach both bounds.
The result sharpens the bounds $1/2$ and $2$ in Skandera--Soskin,
Proposition~11. Known failures of a permanental Koteljanskii inequality
and previously displayed endpoint examples are not claimed as new.
\end{abstract}

\section{Statement and relation to the source question}
A real matrix is \emph{totally positive} if every square minor is strictly
positive, with row and column indices in increasing order. Write
$A_{I,J}$ for its submatrix on $I,J$. Throughout, $12$ denotes $\{1,2\}$,
and likewise for other index strings. Define
\[
 R(A)=\frac{\per(A_{12,12})\per(A_{23,23})}{\per(A)a_{22}}.
\]
All terms in its denominator are positive for a totally positive matrix.
Skandera--Soskin prove $1/2\le R(A)\le2$ in the published
Proposition~11, equation~(49), and state that these bounds are believed
not tight \cite[p.~13]{SS25}. They report observed values $2/3$ and
$121/114$; the latter is an observation, not a conjectured optimal bound.
The sharpness question is posed informally after the proposition, rather than
as a numbered conjecture. The preprint numbers this proposition as 3.6; the FPSAC companion gives
the same bounds and observations \cite{SSfpsac}.
\begin{theorem}\label{main}
For every real totally positive $3\times3$ matrix $A$,
\[
                         \frac23<R(A)<\frac98.
\]
Moreover, for every $c>2/3$ there is such an $A$ with $R(A)<c$, and for
every $c<9/8$ there is such an $A$ with $R(A)>c$. Equivalently,
\[
\begin{split}
 [\,R(A)\ge c\text{ for all totally positive }A\,]&\iff c\le2/3,\\
 [\,R(A)\le c\text{ for all totally positive }A\,]&\iff c\ge9/8.
\end{split}
\]
\end{theorem}
This is a sharp-constant result for the source paper's three-dimensional
ratio. It does not resolve its separate, all-dimensional permanental
Sihong polynomial conjecture, which is the literal question of MathDB
363360.

\section{Normalization and the two certificates}
Multiplying row $i$ by $r_i>0$ and column $j$ by $c_j>0$ preserves total
positivity. Both the numerator and denominator of $R$ acquire the factor
$r_1r_2^2r_3c_1c_2^2c_3$. Thus this scaling preserves $R$.
Taking $r_i=1/a_{i1}$ and $c_j=a_{11}/a_{1j}$ gives
\[
 B=\begin{pmatrix}1&1&1\\1&x&y\\1&z&w\end{pmatrix},\qquad R(B)=R(A).
\]
The $12,12$, $12,23$, and $23,12$ minors imply $x>1$, $y>x$, and $z>x$.
Put $b=y-x>0$, $c=z-x>0$, and
\[
 \Delta=\det B=xw-yz-w+y+z-x>0.
\]
The exact numerator and denominator of the normalized ratio are
\[
 N=(1+x)(xw+yz),\qquad D=x(xw+yz+w+y+z+x)>0.
\]
Expanding both sides gives
\begin{align}
(x-1)(3N-2D)
 &=2x^2(x-1)^2+2x(x+1)(x-1)(b+c)\notag\\
 &\quad +(2x^2+3(x-1))bc+x(x+1)\Delta,\label{lower}\\
(x-1)(9D-8N)
 &=2x^3(x-1)+20x^2(x-1)+2x(x+1)(x-1)(b+c)\notag\\
 &\quad +2(x-2)^2bc+x(x+1)\Delta.\label{upper}
\end{align}
Every summand is nonnegative, and each right-hand side is strictly positive.
Since $x-1>0$ and $D>0$, these identities give $2/3<N/D<9/8$.
This proves both universal bounds for arbitrary, possibly nonsymmetric,
totally positive matrices.

\section{Sharpness inside the strict class}
For completeness, the normalized totally positive class admits a simple
positive parametrization. Given $x>1$ and $b,c,d>0$, set
\[
 y=x+b,\quad z=x+c,\quad w=x+b+c+\frac{bc}{x-1}+d.
\]
Then $\det B=(x-1)d>0$. In lexicographic row-pair and column-pair order,
the nine $2\times2$ minors are
\[
\begin{array}{ccc}
 x-1&x+b-1&b\\
 x+c-1&w-1&w-x-c\\
 c&w-x-b&xw-(x+b)(x+c).
\end{array}
\]
All are positive: in particular the last equals $xd+bc/(x-1)$.
All nine entries are positive too, so this checks every nonempty minor.

For $t>0$, take $x=2$, $b=c=t$, $d=1$. This gives
\[
 U(t)=\begin{pmatrix}1&1&1\\1&2&t+2\\1&t+2&t^2+2t+3\end{pmatrix},
 \qquad R(U(t))=\frac{9t^2+24t+30}{8t^2+24t+38}\longrightarrow\frac98
 \quad(t\longrightarrow\infty).
\]
Every $U(t)$ is strictly totally positive, with determinant $1$.
For $\varepsilon>0$, take $x=1+\varepsilon$, $b=c=d=\varepsilon$, obtaining
\[
 V(\varepsilon)=\begin{pmatrix}
 1&1&1\\1&1+\varepsilon&1+2\varepsilon\\
 1&1+2\varepsilon&1+5\varepsilon
 \end{pmatrix}.
\]
These matrices are strictly totally positive, with determinant
$\varepsilon^2$. Their entries tend to $1$ as $\varepsilon\downarrow0$,
so $R(V(\varepsilon))\to4/6=2/3$. More explicitly,
\[
 R(V(\varepsilon))=
 \frac{(2+\varepsilon)(2+10\varepsilon+9\varepsilon^2)}
 {(1+\varepsilon)(6+20\varepsilon+9\varepsilon^2)}.
\]
These two families prove the optimality assertions in Theorem~\ref{main}.
Strictness of the universal inequalities proves nonattainment.

\section{Prior examples and formal scope}
The failure of $R\le1$ is already known. Skandera--Soskin's
Barrett--Johnson paper gives a totally nonnegative, positive semidefinite
example with $R=144/137$ \cite[\S6]{SSBJ}. Sra's introduction displays
\[
 \begin{pmatrix}1&1&0\\1&2&1\\0&1&1\end{pmatrix},\qquad R=9/8,
\]
so the boundary value $9/8$ itself is also known \cite{Sra}.
These examples, and the failure of $R\le1$, receive no novelty claim here.
Sra's main positive theorem concerns symmetrizable inverse $M$-matrices,
a different class. The proposed contribution of this note is the universal
optimal bounds for \emph{every} totally positive $3\times3$ matrix.
The consulted primary papers do not state these bounds; a bounded
literature search cannot establish absolute priority, and the accompanying independent review records the scope and access limits
of its searches.

The accompanying Lean development contains 53 declarations in six modules,
using Lean~4.34.1 and pinned mathlib revision
\texttt{d13f23b723b8a846827a245b89c10fc7d3f11612}.
It defines the ratio by actual permutation-sum permanents of the two
principal submatrices and the full matrix. Its total-positivity hypothesis
quantifies over every increasing row/column selection of every size at most
three; equivalence with positivity of all nineteen nonempty minors is proved.
The closed endpoint \texttt{source\_sharp\_permanental\_bounds} states both
strict bounds and both optimal-constant equivalences, without normalization
hypotheses. The proof includes positivity of every minor in the sharpness
families and the limiting arguments. No numerical search is used as a
substitute for these formal statements. The accompanying review and replay records document verification of these
statements; the literature comparison remains bounded in scope.

\begin{thebibliography}{9}
\bibitem{SS25} M. Skandera and D. Soskin,
\emph{Permanental inequalities for totally positive matrices},
Electron. J. Combin. \textbf{32}(1) (2025), P1.14.
\href{https://doi.org/10.37236/13098}{doi:10.37236/13098}.
\bibitem{SSfpsac} M. Skandera and D. Soskin,
\emph{Multiplicative permanental inequalities for totally positive matrices},
S\'em. Lothar. Combin. \textbf{93B} (2025), Art. 124.
\bibitem{SSBJ} M. Skandera and D. Soskin,
\emph{Barrett--Johnson inequalities for totally nonnegative matrices},
author manuscript dated May 8, 2026; arXiv:2209.06466.
\bibitem{Sra} S. Sra, \emph{A Koteljanskii inequality for permanents},
arXiv:2609.39979v1 (2026).
\end{thebibliography}
\end{document}
